\section{从写代码到指挥 Agent：工作范式的根本转变}

\subsection{2024年12月：分水岭时刻}

Karpathy 用了一个词来描述自己的状态——\textbf{AI psychosis}（AI 精神错乱）。他说自己一直处于这种状态中，因为个人能力被解锁到了前所未有的程度。过去你的瓶颈是打字速度，而现在有了 Agent，一切都不同了。

\begin{quote}
``Code's not even the right verb anymore. I have to express my will to my agents for 16 hours a day.''

——写代码已经不是对的动词了。我每天得花16个小时向 Agent 表达我的意志。
\end{quote}

他将2024年12月定义为分水岭：在那之前，他自己写80\%的代码、Agent写20\%；之后比例翻转为20/80，而且到现在可能更极端——从12月起他基本没手写过一行代码。

他试图向父母解释这件事，但发现普通人完全没有意识到变化有多剧烈。如果你随便找一个软件工程师看看他的工作台，会发现他做事的默认流程从2024年12月起已经\textbf{完全不同}了。

\begin{importantbox}{核心洞察：约束条件的迁移}
过去至少十年，很多工程任务中人们并不觉得受算力约束。但现在有了这次能力跳跃，你突然发现约束不再是计算资源——\textbf{你自己才是瓶颈}。这其实令人振奋，因为你可以变得更好，所以这件事会上瘾。
\end{importantbox}

Sarah Guo 补充了一个场景：她在 Conviction 基金的工程团队，所有人都不手写代码，工程师们全部佩戴麦克风，对着 Agent 低声说话。她说一开始觉得他们疯了，后来才意识到他们只是走在前面。

\subsection{多 Agent 并行与 Token 吞吐量}

Karpathy 提到 Peter Steinberger 的工作方式作为标杆。Peter 同时在十几个代码仓库之间切换，每个 Agent 用 high effort 模式大约需要20分钟完成任务。操作单位不再是"写一行代码"或"加一个函数"，而是"这个功能交给 Agent 1，那个不会冲突的功能交给 Agent 2"，然后根据你对代码质量的在意程度来审查结果。

\begin{quote}
``What is your token throughput and what token throughput do you command?''

——你的 Token 吞吐量是多少？你能指挥多少 Token 吞吐量？
\end{quote}

他用 PhD 时期的经历做类比：当时 GPU 闲着没跑实验就会焦虑——那代表浪费了算力。现在焦虑的对象从 FLOPS 变成了 Token。如果你还有订阅额度没用完，就意味着你没有最大化自己的 Token 吞吐量。

\begin{knowledgebox}{Peter Steinberger}
奥地利开发者，OpenClaw 开源 Agent 项目的创始人。OpenClaw 是一个可以通过 WhatsApp、Telegram 等消息平台操控的自主 AI Agent，2026年初在 GitHub 上获得超过24万颗星。Steinberger 于2026年2月加入 OpenAI。
\end{knowledgebox}

\subsection{一切都是 Skill Issue}

当被问到现在做项目的瓶颈是什么时，Karpathy 说一切感觉都是 ``skill issue''（技能问题）——不是能力不够，而是你还没找到怎么把现有能力串起来。Agent 的 instructions 写得不够好，memory 工具不够成熟——出问题时总觉得是自己没搞对。

他描述了一个层层递进的结构，像洋葱的层：LLM 被视为理所当然 $\rightarrow$ Agent 被视为理所当然 $\rightarrow$ Claw 实体被视为理所当然 $\rightarrow$ 可以有多个 $\rightarrow$ 可以有 instructions $\rightarrow$ 可以优化 instructions。每一层都无限延伸。

\begin{quote}
``This is why it gets to the psychosis — this is like infinite and everything is skill issue.''

——这就是为什么会到精神错乱的地步——这是无限的，一切都是技能问题。
\end{quote}

\subsection{本章小结}

2024年12月是 AI 编程范式的分水岭。工作单位从代码行变为宏操作（macro actions），衡量指标从 FLOPS 变为 Token 吞吐量，瓶颈从计算资源变为人类自身。多 Agent 并行成为常态，工程师需要发展全新的"肌肉记忆"来适应这种工作方式。


